\documentclass[a4paper,12pt]{ctexart}
\usepackage[left=2.5cm, right=2.5cm, top=2.5cm, bottom=2.5cm]{geometry}
\usepackage{amsmath}
\usepackage{amssymb}
\usepackage{booktabs} %用于三线表
\usepackage{array}
\usepackage{graphicx}
\usepackage{float}
\usepackage{enumitem}
\usepackage{fancyhdr}
\usepackage{xcolor}
\usepackage{tcolorbox} % 用于美化方框

% --- 排版设置 ---
\linespread{1.4} % 行间距
\setlength{\parskip}{0.8em} % 段间距

% --- 页眉页脚 ---
\pagestyle{fancy}
\fancyhf{}
\fancyhead[C]{数值计算方法习题详解}
\fancyfoot[C]{\thepage}

% --- 样式定义宏 ---
% 参数1: 章节标题
\newcommand{\solutionsection}[1]{
    \newpage
    \section*{#1}
    \addcontentsline{toc}{section}{#1}
    \setcounter{equation}{0}
    \hrule height 1pt
    \vspace{1em}
}

% 参数1: 题目编号
% 参数2: 公式与知识点
% 参数3: 解答过程
\newcommand{\solitem}[3]{
    \vspace{1.5em}
    \begin{tcolorbox}[colback=gray!5!white, colframe=gray!50!black, title=\textbf{题目 #1}]
        % 这里留空，题目编号在title中显示
    \end{tcolorbox}
    
    \noindent\textbf{【所用公式与核心考点】}
    \begin{itemize}[leftmargin=2em, label=$\diamond$]
        #2
    \end{itemize}
    
    \noindent\textbf{【详细解答过程】}
    \par
    #3
    \par
    \vspace{1em}
    \hrule
}

\title{\textbf{\Huge 数值计算方法\\ \LARGE 解答汇编 (详细版)}}
\author{整理与解析}
\date{\today}

\begin{document}

\maketitle
\tableofcontents
\newpage

% ==================== 第一章解答 ====================
\solutionsection{第一章：误差与有效数字}

\subsection*{典型例题解答}

\solitem{[3]}{
    \item \textbf{相对误差限定义}：$e_r \approx \frac{\Delta x}{|x^*|}$。
    \item \textbf{有效数字与相对误差的关系定理}：若近似值 $x^*$ 具有 $n$ 位有效数字，则其相对误差限满足：
    \[ |e_r| \le \frac{1}{2a_1} \times 10^{1-n} \]
    其中 $a_1$ 是 $x^*$ 的第一位非零数字（即最高位有效数字）。
}{
    \textbf{第一步：确定首位有效数字}
    设 $\sqrt{20}$ 的近似值为 $x^*$。
    我们需要先估计 $\sqrt{20}$ 的大小范围。
    因为 $\sqrt{16} = 4$ 且 $\sqrt{25} = 5$，所以：
    \[ 4 < \sqrt{20} < 5 \]
    这意味着 $\sqrt{20}$ 的整数部分为 4，写作 $4.\cdots$。
    因此，近似值 $x^*$ 的第一位有效数字为 $a_1 = 4$。
  
    \textbf{第二步：建立不等式}
    题目要求相对误差限 $|e_r| < 0.1\% = 10^{-3}$。
    利用有效数字与相对误差的关系公式：
    \[ |e_r| \le \frac{1}{2 \times a_1} \times 10^{1-n} \]
    我们将 $a_1=4$ 代入，并令其小于目标误差限：
    \[ \frac{1}{2 \times 4} \times 10^{1-n} \le 10^{-3} \]
    \[ \frac{1}{8} \times 10^{1-n} \le 10^{-3} \]
  
    \textbf{第三步：求解有效数字位数 $n$}
    将分数转化为小数便于计算：
    \[ 0.125 \times 10^{1-n} \le 0.001 \]
    两边同时除以 $0.125$：
    \[ 10^{1-n} \le \frac{0.001}{0.125} = \frac{1}{125} = 0.008 \]
    或者直接试算 $n$ 的值：
    \begin{itemize}
        \item 当 $n=3$ 时：$10^{1-3} = 10^{-2} = 0.01$。此时 $0.125 \times 0.01 = 0.00125 > 0.001$，不满足要求。
        \item 当 $n=4$ 时：$10^{1-4} = 10^{-3} = 0.001$。此时 $0.125 \times 0.001 = 0.000125 \le 0.001$，满足要求。
    \end{itemize}
  
    \textbf{结论}：要使 $\sqrt{20}$ 的近似值的相对误差限小于 $0.1\%$，至少要取 \textbf{4} 位有效数字。
}

\solitem{[4]}{
    \item \textbf{积的绝对误差传播公式}：
    对于二元函数 $s = f(l, d) = l \cdot d$，其绝对误差限为：
    \[ \Delta s \approx \left| \frac{\partial s}{\partial l} \right| \Delta l + \left| \frac{\partial s}{\partial d} \right| \Delta d = |d|\Delta l + |l|\Delta d \]
    \item \textbf{积的相对误差传播公式}：
    \[ e_r(ld) \approx e_r(l) + e_r(d) \]
}{
    \textbf{已知条件}：
    矩形长 $l^* = 110\,\mathrm{m}, \Delta l = 0.2\,\mathrm{m}$；
    矩形宽 $d^* = 80\,\mathrm{m}, \Delta d = 0.1\,\mathrm{m}$。
  
    \textbf{1. 计算面积近似值 $s^*$}
    \[ s^* = l^* \cdot d^* = 110 \times 80 = 8800 \, (\mathrm{m}^2) \]
  
    \textbf{2. 计算面积的绝对误差限 $\Delta s$}
    根据误差传播公式：
    \[ \Delta s \approx |d^*| \cdot \Delta l + |l^*| \cdot \Delta d \]
    代入数值：
    \[ \Delta s \approx 80 \times 0.2 + 110 \times 0.1 = 16 + 11 = 27 \, (\mathrm{m}^2) \]
  
    \textbf{3. 计算面积的相对误差限 $e_r(s)$}
    \textit{方法一：定义法}
    \[ e_r(s) = \frac{\Delta s}{s^*} = \frac{27}{8800} \approx 0.003068 \dots \approx 0.31\% \]
  
    \textit{方法二：叠加法}
    先计算长和宽各自的相对误差：
    \[ e_r(l) = \frac{0.2}{110} \approx 0.001818 \]
    \[ e_r(d) = \frac{0.1}{80} = 0.001250 \]
    面积的相对误差为两者之和：
    \[ e_r(s) \approx e_r(l) + e_r(d) = 0.001818 + 0.001250 = 0.003068 \approx 0.31\% \]
  
    \textbf{结论}：面积的绝对误差限为 $27\,\mathrm{m}^2$，相对误差限约为 $0.31\%$。
}

\subsection*{课后习题解答}

\solitem{1}{
    \item \textbf{一元函数误差传播公式}：若 $y = f(x)$，则绝对误差 $\Delta y \approx |f'(x)| \Delta x$。
}{
    设函数为 $y = \ln x$。
    题目已知自变量 $x$ 的相对误差为 $\delta$，即：
    \[ \frac{\Delta x}{x} = \delta \implies \Delta x = x \cdot \delta \]
  
    我们需要计算函数值 $y$ 的绝对误差 $\Delta y$。
    对 $y = \ln x$ 求导得 $y' = \frac{1}{x}$。
    根据误差传播公式：
    \[ \Delta y \approx |y'| \cdot \Delta x = \left| \frac{1}{x} \right| \cdot \Delta x \]
    将 $\Delta x = x \delta$ 代入上式：
    \[ \Delta y \approx \frac{1}{x} \cdot (x \delta) = \delta \]
  
    \textbf{结论}：$\ln x$ 的绝对误差在数值上等于 $x$ 的相对误差 $\delta$。
}

\solitem{2}{
    \item \textbf{幂函数相对误差公式}：若 $y = x^n$，则 $e_r(y) \approx |n| \cdot e_r(x)$。
}{
    已知 $x$ 的相对误差 $e_r(x) = 2\%$。
    目标函数为 $y = x^n$。
  
    推导相对误差：
    \[ e_r(y) = \frac{\Delta y}{y} \approx \frac{|(x^n)'|\Delta x}{x^n} = \frac{|n x^{n-1}| \Delta x}{x^n} = |n| \frac{\Delta x}{x} = |n| e_r(x) \]
    代入已知数据：
    \[ e_r(y) \approx |n| \times 2\% \]
  
    \textbf{结论}：计算 $x^n$ 时，其相对误差约为 $2|n|\%$。
}

\solitem{5}{
    \item \textbf{球体体积公式}：$V = \frac{4}{3}\pi R^3$。
    \item \textbf{反问题误差分析}：已知函数误差限，反求自变量误差限。
}{
    题目要求球体体积 $V$ 的相对误差限不超过 $1\%$，即 $e_r(V) \le 1\%$。
    对于幂函数结构 $V = C \cdot R^3$（其中 $C = \frac{4}{3}\pi$ 为常数，无误差），相对误差传播关系为：
    \[ e_r(V) \approx |3| \cdot e_r(R) = 3 e_r(R) \]
  
    将条件代入：
    \[ 3 e_r(R) \le 1\% \]
    解得半径的相对误差限：
    \[ e_r(R) \le \frac{1\%}{3} \approx 0.33\% \]
  
    \textbf{结论}：为了保证体积精度，度量半径 $R$ 时允许的相对误差限约为 $0.33\%$。
}

\solitem{9}{
    \item \textbf{正方形面积公式}：$S = x^2$。
    \item \textbf{微分估算误差}：$\Delta S \approx S'(x) \Delta x$。
}{
    设正方形边长为 $x \approx 100\,\mathrm{cm}$。
    面积函数 $S(x) = x^2$。
    题目要求面积的绝对误差 $\Delta S \le 1\,\mathrm{cm}^2$。
  
    利用微分公式建立关系：
    \[ \Delta S \approx 2x \cdot \Delta x \]
    代入数值：
    \[ 2 \times 100 \times \Delta x \le 1 \]
    \[ 200 \Delta x \le 1 \]
    解得边长误差限：
    \[ \Delta x \le \frac{1}{200} = 0.005\,\mathrm{cm} \]
  
    \textbf{结论}：测量边长时，误差不应超过 $0.005\,\mathrm{cm}$（这意味着测量工具需精确到小数点后2位以上）。
}

\newpage

% ==================== 第二章解答 ====================
\solutionsection{第二章：插值法}

\subsection*{典型例题解答}

\solitem{[2]}{
    \item \textbf{拉格朗日线性插值}：$L_1(x) = y_0 \frac{x-x_1}{x_0-x_1} + y_1 \frac{x-x_0}{x_1-x_0}$。
    \item \textbf{拉格朗日抛物插值}：$L_2(x) = \sum_{i=0}^2 y_i l_i(x)$。
    \item \textbf{截断误差余项}：
    $R_1(x) = \frac{f''(\xi)}{2!}(x-x_0)(x-x_1)$；
    $R_2(x) = \frac{f'''(\xi)}{3!}(x-x_0)(x-x_1)(x-x_2)$。
}{
    已知数据点：
    $x_0 = 0.32, y_0 = \sin(0.32) = 0.314567$；
    $x_1 = 0.34, y_1 = \sin(0.34) = 0.333487$；
    $x_2 = 0.36, y_2 = \sin(0.36) = 0.352274$。
    待求点 $x = 0.3367$。
  
    \textbf{1. 线性插值 (Linear Interpolation)}
    选取距离 $x$ 最近的两个节点 $x_0, x_1$（因为 $0.32 < 0.3367 < 0.34$）。
    插值公式：
    \[ L_1(0.3367) = 0.314567 + \frac{0.333487 - 0.314567}{0.34 - 0.32}(0.3367 - 0.32) \]
    计算斜率：
    \[ \frac{\Delta y}{\Delta x} = \frac{0.018920}{0.02} = 0.946 \]
    计算增量：
    \[ L_1 = 0.314567 + 0.946 \times 0.0167 = 0.314567 + 0.0157982 = 0.3303652 \]
  
    \textbf{误差估计}：$f(x) = \sin x \Rightarrow |f''(x)| = |-\sin x|$。在 $[0.32, 0.34]$ 上最大值约为 $\sin(0.34) \approx 0.3335$。
    \[ |R_1| \le \frac{0.3335}{2} |(0.0167)(-0.0033)| \approx 0.16675 \times 0.00005511 \approx 0.9 \times 10^{-5} \]
  
    \textbf{2. 抛物插值 (Quadratic Interpolation)}
    利用 $x_0, x_1, x_2$ 三点构造二次多项式。
    基函数计算：
    \begin{align*}
    l_0(0.3367) &= \frac{(0.3367-0.34)(0.3367-0.36)}{(0.32-0.34)(0.32-0.36)} = \frac{(-0.0033)(-0.0233)}{(-0.02)(-0.04)} = \frac{0.00007689}{0.0008} \approx 0.09611 \\
    l_1(0.3367) &= \frac{(0.3367-0.32)(0.3367-0.36)}{(0.34-0.32)(0.34-0.36)} = \frac{(0.0167)(-0.0233)}{(0.02)(-0.02)} = \frac{-0.00038911}{-0.0004} \approx 0.97278 \\
    l_2(0.3367) &= \frac{(0.3367-0.32)(0.3367-0.34)}{(0.36-0.32)(0.36-0.34)} = \frac{(0.0167)(-0.0033)}{(0.04)(0.02)} = \frac{-0.00005511}{0.0008} \approx -0.06889
    \end{align*}
    组合结果：
    \[ L_2 \approx 0.314567(0.09611) + 0.333487(0.97278) + 0.352274(-0.06889) \]
    \[ L_2 \approx 0.030233 + 0.324410 - 0.024268 = 0.330375 \]
  
    \textbf{误差估计}：$|f'''(x)| = |-\cos x| \le |\cos 0.32| \approx 0.949$。
    \[ |R_2| \le \frac{0.949}{6} |(0.0167)(-0.0033)(-0.0233)| \approx 0.158 \times 1.28 \times 10^{-6} \approx 2.0 \times 10^{-7} \]
  
    \textbf{结论}：线性插值结果 $0.330365$，误差量级 $10^{-5}$；抛物插值结果 $0.330375$，误差量级 $10^{-7}$。
}

\solitem{[4]}{
    \item \textbf{牛顿插值多项式}：
    $P_n(x) = f(x_0) + f[x_0,x_1](x-x_0) + \cdots + f[x_0,\dots,x_n]\prod_{i=0}^{n-1}(x-x_i)$。
    \item \textbf{节点选择策略}：插值点应尽量靠近目标点 $x$，以减小误差积项 $\prod (x-x_i)$ 的值。
}{
    \textbf{1. 节点选择与排序}
    目标点 $x=0.596$。为了保证精度，选取距离 $0.596$ 最近的 5 个点进行 4 次插值。
    距离排序：$0.55, 0.65, 0.40, 0.80, 0.90$。
    重排节点顺序为：$x_0=0.55, x_1=0.65, x_2=0.40, x_3=0.80, x_4=0.90$。
  
    \textbf{2. 构造差商表}
    \begin{table}[H]
    \centering
    \small
    \begin{tabular}{c c l l l l}
    \toprule
    $x_i$ & $f(x_i)$ & 一阶差商 & 二阶差商 & 三阶差商 & 四阶差商 \\ 
    \midrule
    0.55 & 0.57815 & & & & \\
         &         & 1.18600 & & & \\
    0.65 & 0.69675 &         & 1.13500 & & \\
         &         & 1.14400 &         & 1.29200 & \\
    0.40 & 0.41075 &         & 1.26422 &         & 1.63600 \\
         &         & 1.19340 &         & 0.88286 & \\
    0.80 & 0.88811 &         & 1.57320 &         & \\
         &         & 1.38410 &         & & \\
    0.90 & 1.02652 &         & & & \\
    \bottomrule
    \end{tabular}
    \end{table}
    \textit{注：计算示例 $f[x_0, x_1] = \frac{0.69675-0.57815}{0.65-0.55} = 1.186$。}
  
    \textbf{3. 构建牛顿插值多项式}
    取差商表的主对角线元素作为系数：
    \begin{align*}
    P_4(x) &= 0.57815 \\
           &+ 1.186(x-0.55) \\
           &+ 1.135(x-0.55)(x-0.65) \\
           &+ 1.292(x-0.55)(x-0.65)(x-0.40) \\
           &+ 1.636(x-0.55)(x-0.65)(x-0.40)(x-0.80)
    \end{align*}
  
    \textbf{4. 代入求解}
    令 $x = 0.596$，计算各乘积项：
    $t_0 = x-0.55 = 0.046$
    $t_1 = x-0.65 = -0.054$
    $t_2 = x-0.40 = 0.196$
    $t_3 = x-0.80 = -0.204$
  
    \begin{align*}
    P_4(0.596) &= 0.57815 \\
               &+ 1.186(0.046) \quad (\approx 0.054556) \\
               &+ 1.135(0.046)(-0.054) \quad (\approx -0.002819) \\
               &+ 1.292(0.046)(-0.054)(0.196) \quad (\approx -0.000629) \\
               &+ 1.636(0.046)(-0.054)(0.196)(-0.204) \quad (\approx 0.000163)
    \end{align*}
    求和：
    \[ 0.57815 + 0.054556 - 0.002819 - 0.000629 + 0.000163 = 0.629421 \]
  
    \textbf{结果}：$f(0.596) \approx 0.62942$。
}

\subsection*{课后习题解答}

\solitem{1}{
    \item \textbf{插值多项式的唯一性}：无论采用何种基底（单项式、拉格朗日、牛顿），只要满足插值条件，化简后的多项式是唯一的。
}{
    数据点：$(-1, -3), (1, 0), (2, 4)$。
  
    \textbf{方法一：单项式基底法 (待定系数)}
    设 $P(x) = a + bx + cx^2$。代入三个点得到线性方程组：
    \[ \begin{cases} a - b + c = -3 \quad (1) \\ a + b + c = 0 \quad \quad (2) \\ a + 2b + 4c = 4 \quad (3) \end{cases} \]
    求解过程：
    (2)-(1) 得 $2b = 3 \Rightarrow b = 1.5$。
    将 $b$ 代入(2) 得 $a+c = -1.5$。
    将 $b$ 代入(3) 得 $a+4c = 4 - 3 = 1$。
    两式相减 $3c = 2.5 \Rightarrow c = 5/6$。
    $a = -1.5 - 5/6 = -14/6 = -7/3$。
    \[ P(x) = -\frac{7}{3} + \frac{3}{2}x + \frac{5}{6}x^2 \]
  
    \textbf{方法二：拉格朗日插值法}
    \[ L_2(x) = (-3)\frac{(x-1)(x-2)}{(-1-1)(-1-2)} + 0 + 4\frac{(x+1)(x-1)}{(2+1)(2-1)} \]
    \[ = (-3)\frac{x^2-3x+2}{6} + 4\frac{x^2-1}{3} \]
    \[ = -\frac{1}{2}(x^2-3x+2) + \frac{4}{3}(x^2-1) \]
    合并同类项：
    $x^2$项：$-\frac{1}{2} + \frac{4}{3} = \frac{5}{6}$
    $x$项：$\frac{3}{2}$
    常数项：$-1 - \frac{4}{3} = -\frac{7}{3}$
    \[ P(x) = \frac{5}{6}x^2 + \frac{3}{2}x - \frac{7}{3} \]
  
    \textbf{方法三：牛顿插值法}
    一阶差商：$\frac{0-(-3)}{1-(-1)} = 1.5$, $\frac{4-0}{2-1} = 4$
    二阶差商：$\frac{4-1.5}{2-(-1)} = \frac{2.5}{3} = \frac{5}{6}$
  
    $N_2(x) = -3 + 1.5(x+1) + \frac{5}{6}(x+1)(x-1)$
    $= -3 + 1.5x + 1.5 + \frac{5}{6}(x^2-1)$
    $= \frac{5}{6}x^2 + 1.5x - 1.5 - \frac{5}{6} = \frac{5}{6}x^2 + \frac{3}{2}x - \frac{7}{3}$
  
    \textbf{结论}：三种方法所得多项式完全一致。
}

\newpage

% ==================== 第三章解答 ====================
\solutionsection{第三章：曲线拟合与逼近}

\subsection*{典型例题解答}

\solitem{[9]}{
    \item \textbf{加权最小二乘法原理}：
    寻找 $y = a + bx$ 使得目标函数 $E = \sum_{i=1}^n \omega_i (y_i - (a+bx_i))^2$ 最小。
    \item \textbf{法方程组}：
    \[ \begin{pmatrix} \sum \omega_i & \sum \omega_i x_i \\ \sum \omega_i x_i & \sum \omega_i x_i^2 \end{pmatrix} \begin{pmatrix} a \\ b \end{pmatrix} = \begin{pmatrix} \sum \omega_i y_i \\ \sum \omega_i x_i y_i \end{pmatrix} \]
}{
    \textbf{1. 列表计算各项加权和}
    \begin{table}[H]
    \centering
    \begin{tabular}{c c c c c c c c}
    \toprule
    $i$ & $x_i$ & $y_i$ & $\omega_i$ & $\omega_i x_i$ & $\omega_i x_i^2$ & $\omega_i y_i$ & $\omega_i x_i y_i$ \\ 
    \midrule
    1 & 1 & 4 & 2 & 2 & 2 & 8 & 8 \\
    2 & 2 & 4.5 & 1 & 2 & 4 & 4.5 & 9 \\
    3 & 3 & 6 & 3 & 9 & 27 & 18 & 54 \\
    4 & 4 & 8 & 1 & 4 & 16 & 8 & 32 \\
    5 & 5 & 8.5 & 1 & 5 & 25 & 8.5 & 42.5 \\ 
    \midrule
    $\sum$ & & & \textbf{8} & \textbf{22} & \textbf{74} & \textbf{47} & \textbf{145.5} \\
    \bottomrule
    \end{tabular}
    \end{table}
  
    \textbf{2. 建立法方程组}
    根据公式代入求和值：
    \[ \begin{cases} 8a + 22b = 47 \quad (1) \\ 22a + 74b = 145.5 \quad (2) \end{cases} \]
  
    \textbf{3. 求解线性方程组}
    由(1)得 $a = (47 - 22b)/8$。代入(2)：
    \[ 22\left(\frac{47 - 22b}{8}\right) + 74b = 145.5 \]
    \[ 2.75(47 - 22b) + 74b = 145.5 \]
    \[ 129.25 - 60.5b + 74b = 145.5 \]
    \[ 13.5b = 16.25 \Rightarrow b = \frac{16.25}{13.5} = \frac{1625}{1350} = \frac{65}{54} \approx 1.2037 \]
    回代求 $a$：
    \[ a = \frac{47 - 22(1.2037)}{8} = \frac{47 - 26.4814}{8} = 2.5648 \]
  
    \textbf{结论}：拟合直线方程为 $y = 2.5648 + 1.2037x$。
}

\solitem{[10]}{
    \item \textbf{非线性拟合线性化}：
    对于 $y = ae^{bx}$，两边取自然对数得 $\ln y = \ln a + bx$。
    令 $Y = \ln y, A = \ln a, B = b$，转化为 $Y = A + Bx$ 的线性最小二乘问题。
}{
    \textbf{1. 数据预处理与转化}
    $x_i$: 1.00, 1.25, 1.50, 1.75, 2.00
    $y_i$: 5.10, 5.79, 6.53, 7.45, 8.46
    计算 $Y_i = \ln y_i$: 1.6292, 1.7561, 1.8764, 2.0082, 2.1353
  
    \textbf{2. 计算法方程系数 ($n=5$)}
    $\sum x = 7.5$
    $\sum x^2 = 1 + 1.5625 + 2.25 + 3.0625 + 4 = 11.875$
    $\sum Y = 1.6292 + \dots + 2.1353 = 9.4052$
    $\sum xY = 1(1.6292) + 1.25(1.7561) + 1.5(1.8764) + 1.75(2.0082) + 2(2.1353)$
    $= 1.6292 + 2.1951 + 2.8146 + 3.5144 + 4.2706 = 14.4239$
  
    \textbf{3. 求解 $A$ 和 $B$}
    法方程组：
    \[ \begin{cases} 5A + 7.5B = 9.4052 \\ 7.5A + 11.875B = 14.4239 \end{cases} \]
    利用克拉默法则或消元法：
    系数行列式 $D = 5 \times 11.875 - 7.5^2 = 59.375 - 56.25 = 3.125$
    $D_A = 9.4052 \times 11.875 - 7.5 \times 14.4239 = 3.5075$
    $D_B = 5 \times 14.4239 - 7.5 \times 9.4052 = 1.5805$
  
    $A = D_A / D = 1.1224$
    $B = D_B / D = 0.50576$
  
    \textbf{4. 还原参数}
    $b = B = 0.5058$
    $a = e^A = e^{1.1224} \approx 3.072$
  
    \textbf{结论}：拟合曲线为 $y = 3.072 e^{0.5058x}$。
}

\newpage

% ==================== 第四章解答 ====================
\solutionsection{第四章：数值积分与微分}

\subsection*{典型例题解答}

\solitem{[1]}{
    \item \textbf{代数精度定义}：若求积公式对 $f(x) = 1, x, x^2, \dots, x^m$ 均精确成立，而对 $x^{m+1}$ 不精确，则称该公式具有 $m$ 阶代数精度。
    \item \textbf{待定系数法}：通过令公式对幂函数精确成立来求解参数。
}{
    设求积公式结构为 $\int_0^1 f(x)dx \approx A_0 f(0) + A_1 f(1) + B_0 f'(0)$。
    我们要确定 3 个参数，期望达到至少 2 阶代数精度。令 $f(x)$ 依次取 $1, x, x^2$。
  
    \textbf{1. 建立方程组}
    \begin{itemize}
        \item 取 $f(x) = 1$，则 $\int_0^1 1 dx = 1$。右端 $= A_0 + A_1$。
        \item 取 $f(x) = x$，则 $\int_0^1 x dx = 1/2$。$f'(x)=1 \Rightarrow f'(0)=1$。右端 $= A_1 + B_0$。
        \item 取 $f(x) = x^2$，则 $\int_0^1 x^2 dx = 1/3$。$f'(x)=2x \Rightarrow f'(0)=0$。右端 $= A_1$。
    \end{itemize}
    得到方程组：
    \[ \begin{cases} A_0 + A_1 = 1 \\ A_1 + B_0 = 1/2 \\ A_1 = 1/3 \end{cases} \]
  
    \textbf{2. 求解参数}
    由第三式直接得 $A_1 = 1/3$。
    代入第二式：$1/3 + B_0 = 1/2 \Rightarrow B_0 = 1/6$。
    代入第一式：$A_0 + 1/3 = 1 \Rightarrow A_0 = 2/3$。
  
    \textbf{3. 检查代数精度}
    测试 $f(x) = x^3$：
    左边 $= \int_0^1 x^3 dx = 1/4$。
    右边 $= \frac{2}{3}(0)^3 + \frac{1}{3}(1)^3 + \frac{1}{6}(3 \cdot 0^2) = 1/3$。
    因为 $1/4 \neq 1/3$，所以不具有 3 阶精度。
  
    \textbf{结论}：$A_0 = 2/3, A_1 = 1/3, B_0 = 1/6$。公式具有 2 阶代数精度。
}

\solitem{[3]}{
    \item \textbf{复合梯形公式}：$T_n = \frac{h}{2} [f(a) + 2\sum_{i=1}^{n-1} f(x_i) + f(b)]$。
    \item \textbf{复合辛普森公式}：$S_n = \frac{h}{3} [f(a) + 4\sum_{odd} + 2\sum_{even} + f(b)]$。
    \item \textbf{龙格误差估计}：$I - T_{2n} \approx \frac{1}{3}(T_{2n} - T_n)$。
}{
    积分区间 $[0, 1]$，分点数 $n=8$，步长 $h = 0.125$。
    数据点：$x_i = 0, 0.125, 0.25, \dots, 1$。
    计算得：
    $f(0)=1, f(1)=0.8415$。
    $\sum f(x_i)$ (中间点) $= 6.6448$。
  
    \textbf{1. 计算复合梯形值 $T_8$}
    \[ T_8 = \frac{0.125}{2} [1 + 0.8415 + 2 \times 6.6448] = 0.0625 [15.1311] = 0.945691 \]
  
    \textbf{2. 计算复合辛普森值 $S_8$}
    奇数点和 $\Sigma_{odd} \approx 3.7875$；偶数点和 $\Sigma_{even} \approx 2.8573$。
    \[ S_8 = \frac{0.125}{3} [1.8415 + 4(3.7875) + 2(2.8573)] \approx 0.946083 \]
  
    \textbf{3. 误差估计 (利用 $T_4, S_4$)}
    计算得 $T_4 \approx 0.944514$。
    $T_8$ 的误差估计：
    \[ R \approx \frac{1}{3}(T_8 - T_4) = \frac{1}{3}(0.945691 - 0.944514) \approx 3.9 \times 10^{-4} \]
}

\newpage

% ==================== 第五章解答 ====================
\solutionsection{第五章：线性方程组的直接解法}

\subsection*{典型例题解答}

\solitem{[5]}{
    \item \textbf{Doolittle 分解 (LU分解)}：将矩阵 $A$ 分解为单位下三角矩阵 $L$ 和上三角矩阵 $U$。
    \item \textbf{求解步骤}：
    1. 分解 $A=LU$。
    2. 解 $Ly=b$ (前代法)。
    3. 解 $Ux=y$ (回代法)。
}{
    方程组 $Ax=b$，$A = \begin{pmatrix} 1 & 2 & 3 \\ 2 & 5 & 2 \\ 3 & 1 & 5 \end{pmatrix}, b = \begin{pmatrix} 14 \\ 18 \\ 20 \end{pmatrix}$。
  
    \textbf{1. LU 分解过程}
    \textit{第1步} ($U$第1行, $L$第1列):
    $u_{11}=1, u_{12}=2, u_{13}=3$。
    $l_{11}=1, l_{21}=2/1=2, l_{31}=3/1=3$。
  
    \textit{第2步} ($U$第2行, $L$第2列):
    $u_{22} = a_{22} - l_{21}u_{12} = 5 - 2\times 2 = 1$。
    $u_{23} = a_{23} - l_{21}u_{13} = 2 - 2\times 3 = -4$。
    $l_{32} = (a_{32} - l_{31}u_{12})/u_{22} = (1 - 3\times 2)/1 = -5$。
  
    \textit{第3步} ($U$第3行):
    $u_{33} = a_{33} - (l_{31}u_{13} + l_{32}u_{23}) = 5 - (3\times 3 + (-5)\times(-4)) = 5 - 29 = -24$。
  
    \[ L = \begin{pmatrix} 1 & 0 & 0 \\ 2 & 1 & 0 \\ 3 & -5 & 1 \end{pmatrix}, \quad U = \begin{pmatrix} 1 & 2 & 3 \\ 0 & 1 & -4 \\ 0 & 0 & -24 \end{pmatrix} \]
  
    \textbf{2. 求解 Ly = b}
    \[ \begin{cases} y_1 = 14 \\ 2y_1 + y_2 = 18 \Rightarrow y_2 = 18 - 28 = -10 \\ 3y_1 - 5y_2 + y_3 = 20 \Rightarrow 42 + 50 + y_3 = 20 \Rightarrow y_3 = -72 \end{cases} \]
  
    \textbf{3. 求解 Ux = y}
    \[ \begin{cases} -24x_3 = -72 \Rightarrow x_3 = 3 \\ x_2 - 4x_3 = -10 \Rightarrow x_2 - 12 = -10 \Rightarrow x_2 = 2 \\ x_1 + 2x_2 + 3x_3 = 14 \Rightarrow x_1 + 4 + 9 = 14 \Rightarrow x_1 = 1 \end{cases} \]
  
    \textbf{结论}：方程组的解为 $x = (1, 2, 3)^T$。
}

\newpage

% ==================== 第七章解答 ====================
\solutionsection{第七章：非线性方程求根}

\subsection*{典型例题解答}

\solitem{[7]}{
    \item \textbf{牛顿迭代公式}：$x_{k+1} = x_k - \frac{f(x_k)}{f'(x_k)}$。
    \item \textbf{收敛性}：当初值 $x_0$ 充分接近单根 $\xi$ 时，牛顿法具有平方收敛速度。
}{
    求解方程 $xe^x - 1 = 0$。
    设 $f(x) = xe^x - 1$，则 $f'(x) = e^x + xe^x = e^x(1+x)$。
  
    \textbf{1. 构造迭代格式}
    \[ x_{k+1} = x_k - \frac{x_k e^{x_k} - 1}{e^{x_k}(1+x_k)} \]
  
    \textbf{2. 确定初始值}
    $f(0) = -1 < 0, f(1) = e - 1 > 0$，根在 $(0, 1)$ 内。取 $x_0 = 0.5$。
  
    \textbf{3. 迭代计算}
    \textit{第1次迭代}：
    $f(0.5) = 0.5e^{0.5} - 1 \approx -0.1756$
    $f'(0.5) = e^{0.5}(1.5) \approx 2.4731$
    $x_1 = 0.5 - \frac{-0.1756}{2.4731} \approx 0.5710$
  
    \textit{第2次迭代}：
    $f(0.5710) \approx 0.0107$
    $f'(0.5710) \approx 2.780$
    $x_2 = 0.5710 - \frac{0.0107}{2.780} \approx 0.5672$
  
    \textit{第3次迭代}：
    $x_3 \approx 0.5671$。此时 $f(x_3)$ 已极小。
  
    \textbf{结论}：方程的根约为 $x^* \approx 0.5671$。
}

\end{document}